\section{Discussion}
\label{sec:discussion}

The pilot supports syntactic portability but not yet behavioral conformance: one
unchanged structured profile was accepted by two heterogeneous model families,
yet the small automatic analysis cannot show that either family produced answers
the target user would prefer. That distinction prevents a successful serializer
or prompt injector from being mislabeled as successful personalization.

The compiler ablation also exposes a cost-accounting trap. A compressed policy can
save input tokens while increasing generated tokens and wall-clock latency. In our
sample, directive-like compact clauses appear to have encouraged fuller responses
despite a global concise preference. This interpretation is mechanistic and
post-hoc; output caps and unequal model sizes are alternative explanations. Future
work should jointly optimize adherence, correctness, input tokens, output tokens,
and truncation rather than treating prompt length as the efficiency endpoint.

The near-zero compact E-versus-history difference on reference coverage should not
be read as equivalence. The interval is wide, the diagnostic is lexical, and only
one synthetic persona and four tasks were used. Conversely, the absence of a
positive automatic result is useful counterevidence against assuming that typed
profiles automatically improve answers. The condition-blinded files now make the
next decisive step concrete: target-user preference judgments and independent
correctness review under a frozen analysis plan.

The larger research program should distinguish broad benefit, style-only benefit,
provider and domain heterogeneity, and net harm from incorrect or stale preferences.
Practical effect sizes, correction burden, and abstention quality matter more than
statistical significance alone.

Implementation of contextual negative feedback further sharpens the research
program. Asking whether a stored preference is true and asking whether it belongs
in a current prompt are separable labeling tasks. Future human evaluation should
therefore report inference calibration, applicability precision, abstention
coverage, and recovery after each error type rather than collapsing them into one
personalization score.
